
\section{Next-Game Prediction}

The statistical analysis of historical match records
provides a basis for estimating the scoring performance
of the next game. This section develops a prediction
model using the recorded scores from previous matches.

The objective is to estimate the points likely to be
scored by our team and the opponent in the next game.
The prediction is performed independently for Singles,
Doubles, and Mixed Doubles.

% ==================================================
\subsection{Prediction Objective}
% ==================================================

The prediction model aims to estimate three quantities:

\begin{enumerate}
    \item Expected points scored by our team.
    \item Expected points scored by the opponent.
    \item Expected net result of the next game.
\end{enumerate}

The predicted net result provides an indication of
the expected scoring advantage.

A positive predicted net result suggests that our
team may score more points, while a negative value
suggests that the opponent may score more points.

These predictions are estimates based on historical
data and do not guarantee the outcome of a future game.

% ==================================================
\subsection{Least-Squares Linear Regression}
% ==================================================

Least-squares linear regression is used to model the
relationship between the chronological match index
and the recorded scores.

Let the historical matches be represented by

\[
(x_i,W_i), \qquad i=1,2,\ldots,n
\]

where:

\begin{itemize}
    \item $x_i$ is the chronological index of the
    $i$-th recorded match.
    \item $W_i$ is the points scored by our team.
    \item $n$ is the number of valid match records.
\end{itemize}

A linear model is expressed as

\[
\widehat{W}(x)=a_W+b_Wx
\]

where $a_W$ is the intercept and $b_W$ is the slope.

The slope and intercept are determined by minimizing
the sum of squared prediction errors:

\[
S_W =
\sum_{i=1}^{n}
\left[W_i-(a_W+b_Wx_i)\right]^2
\]

The least-squares estimates are

\[
b_W =
\frac{
n\sum x_iW_i-(\sum x_i)(\sum W_i)
}{
n\sum x_i^2-(\sum x_i)^2
}
\]

and

\[
a_W=\overline{W}-b_W\overline{x}
\]

The same procedure is applied independently to the
opponent's scores.

Thus, the opponent's scoring model is

\[
\widehat{L}(x)=a_L+b_Lx
\]

where $L_i$ represents the points scored by the
opponent.

% ==================================================
\subsection{Predicting the Next Game}
% ==================================================

If $n$ valid records are available for a particular
match format, the next chronological index is

\[
x_{\text{next}}=n+1
\]

The predicted points scored by our team are

\[
\widehat{W}_{\text{next}}
=
a_W+b_W(n+1)
\]

Similarly, the predicted points scored by the
opponent are

\[
\widehat{L}_{\text{next}}
=
a_L+b_L(n+1)
\]

The predicted net result is then calculated as

\[
\widehat{R}_{\text{next}}
=
\widehat{W}_{\text{next}}
-
\widehat{L}_{\text{next}}
\]

This gives an estimated scoring comparison for
the next game.

% ==================================================
\subsection{Separate Models for Match Formats}
% ==================================================

The prediction process is performed independently
for all three match formats.

\subsubsection{Singles Prediction}

The Singles model uses only historical Singles
records. Separate regression equations are fitted
for our team's points and the opponent's points.

The resulting equations are used to estimate the
next Singles score.

\subsubsection{Doubles Prediction}

The Doubles model uses only historical Doubles
records. This ensures that the prediction is based
on the scoring pattern observed in Doubles matches.

Separate equations are fitted for both teams to
estimate the next Doubles score.

\subsubsection{Mixed Doubles Prediction}

The Mixed Doubles model uses only historical Mixed
Doubles records. The same least-squares procedure
is applied to estimate the next Mixed Doubles score.

Using separate models avoids combining scoring
patterns from different match formats.

% ==================================================
\subsection{Prediction Summary}
% ==================================================

The final predictions can be presented in the
following table after evaluating the regression
equations.

\begin{table}[H]
\centering
\caption{Estimated Scores for the Next Game}
\begin{tabular}{lccc}
\toprule
\textbf{Format} &
\textbf{Our Points} &
\textbf{Opponent Points} &
\textbf{Net Result} \\
\midrule
Singles & -- & -- & -- \\
Doubles & -- & -- & -- \\
Mixed Doubles & -- & -- & -- \\
\bottomrule
\end{tabular}
\end{table}

The predicted scores are continuous numerical
estimates. Since actual match scores are integers,
the estimates may be rounded to the nearest
non-negative integer for presentation.

% ==================================================
\subsection{Limitations of the Model}
% ==================================================

The prediction model is based on historical scoring
records and assumes that the observed linear trend
provides useful information about future performance.

However, several factors can influence the actual
result:

\begin{itemize}
    \item Changes in the players participating in
    a match.
    \item Differences in opponent strength.
    \item Player fatigue and physical condition.
    \item Variations in playing conditions.
    \item The relatively small number of records
    available for some match formats.
\end{itemize}

The model does not explicitly account for these
factors. Therefore, its output should be interpreted
as a statistical estimate rather than a guaranteed
match result.

% ==================================================
\subsection{Conclusion}
% ==================================================

Least-squares linear regression provides a simple
method for estimating the next game's scoring
performance from historical records.

By fitting separate models for Singles, Doubles,
and Mixed Doubles, the prediction process preserves
the differences between the three match formats.

The predicted scores and net results provide a
quantitative estimate of future scoring performance
and complement the descriptive statistical analysis
presented in the previous section.

